% !TeX root = ../presentacion.tex
% ===========================================================================
%  3. Arquitectura — 0:50
%
%  Bloque bisagra: no puntúa por sí solo, pero hace inteligibles los dos
%  siguientes (15 % + 15 %).
%
%  POR QUÉ UN MAPA PROPIO Y NO LA VISTA C4 EXPORTADA
%  -------------------------------------------------
%  La vista de contenedores (V-02) mide 2254 x 2720 puntos: es más alta que
%  ancha. Proyectada en una lámina 16:9 se queda en un quinto del ancho y
%  sus etiquetas dejan de leerse desde el fondo del aula. Este mapa dice lo
%  mismo con los mismos colores de Structurizr y sí se lee.
%
%  Las cuatro vistas C4 completas están en el respaldo, a lámina entera,
%  para proyectarlas si alguien pregunta por el detalle.
% ===========================================================================

\bloque{Arquitectura}{0:50}{\exponeArquitectura}%
       {Base de los criterios de microfrontends y microservicios (30\,\%)}

\begin{frame}{Once contenedores, un solo origen}
  \note{Recorrer el mapa de arriba abajo en veinte segundos. El detalle que
  merece énfasis: el edge es el único conectado a las dos redes de Docker,
  y por eso el navegador no puede alcanzar un microservicio saltándose el
  gateway. El diagrama C4 completo está en el respaldo.}

  \begin{center}
  \setlength{\tabcolsep}{3pt}
  \begin{tabular}{@{}*{4}{p{0.232\textwidth}}@{}}
    \multicolumn{4}{@{}p{0.98\textwidth}@{}}{%
      \caja{tintaEdge}{Navegador $\rightarrow$ \id{edge} (nginx)}{un único origen: sirve el shell, los tres remotes y \id{/api}}}\\[0.5pt]
    \multicolumn{4}{@{}c@{}}{\bajada{Module Federation: los remotes se resuelven en ejecución}}\\[0.5pt]
    \caja{tintaFront}{\id{shell}}{host} &
    \caja{tintaFront}{\id{mf-reservas}}{remote} &
    \caja{tintaFront}{\id{mf-adminis\-tracion}}{remote} &
    \caja{tintaFront}{\id{mf-reportes}}{remote} \\[0.5pt]
    \multicolumn{4}{@{}c@{}}{\bajada{\id{/api} — entrada única}}\\[0.5pt]
    \multicolumn{4}{@{}p{0.98\textwidth}@{}}{%
      \caja{tintaGw}{\id{api-gateway}}{autentica y propaga \id{X-User-Id} y \id{X-User-Role}}}\\[0.5pt]
    \multicolumn{4}{@{}c@{}}{\bajada{HTTP/REST síncrono}}\\[0.5pt]
    \caja{tintaMs}{\id{ms-usuarios}}{} &
    \caja{tintaMs}{\id{ms-canchas}}{} &
    \caja{tintaMs}{\id{ms-reservas}}{} &
    \caja{tintaMs}{\id{ms-reportes}}{sin base: agrega por REST} \\[0.5pt]
    \caja{tintaDb}{\id{usuarios\_db}}{} &
    \caja{tintaDb}{\id{canchas\_db}}{} &
    \caja{tintaDb}{\id{reservas\_db}}{} &
    \\
  \end{tabular}
  \end{center}

  \vspace{-6pt}
  {\tiny\color{textoSuave} Modelo C4 en Structurizr DSL. Las cuatro vistas
  completas, en el respaldo.}
\end{frame}
